\documentclass[12pt]{article}
\usepackage{graphicx} % Required for inserting images
\usepackage[margin=0.5in]{geometry}
\usepackage{amsmath, amssymb}
\usepackage{mathrsfs}


\usepackage{soul}


% Dr. David Brown Commands
\newcommand{\ds}{\displaystyle}
\newcommand{\R}{\mathbb{R}}
\newcommand{\N}{\mathbb{N}}
\newcommand{\Q}{\mathbb{Q}}
\newcommand{\C}{\mathbb{C}}


% Commands to get script r
\usepackage{calligra}
\DeclareMathAlphabet{\mathcalligra}{T1}{calligra}{m}{n}
\DeclareFontShape{T1}{calligra}{m}{n}{<->s*[2.2]callig15}{}
\newcommand{\scripty}[1]{\ensuremath{\mathcalligra{#1}}}

% Unit commands
\newcommand{\degree}{\text{°}}
\newcommand{\unitSpeed}{\frac{\text{m}}{\text{s}}}
\newcommand{\unitAcceleration}{\frac{\text{m}}{\text{s}^2}}

% Constant commands
\newcommand{\avogadro}{6.022\times10^{23}}

\newcommand{\boltzmannConstK}{1.381\times10^{-23}\frac{\text{J}}{\text{K}}}
\newcommand{\boltzmannConstInEV}{8.617\times10^{-5}\frac{\text{eV}}{\text{K}}}
\newcommand{\planckConst}{6.626\times10^{-34}\text{J}\cdot\text{s}}

\newcommand{\atmP}{1.01325\times10^5\,\text{Pa}}

% Known Value Commands
\newcommand{\electronMass}{9.109\times10^{-31}\,\text{kg}}
\newcommand{\protonMass}{1.673\times10^{-27}\,\text{kg}}

% Commands to easily write partial derivatives
\newcommand{\partDeriv}[2]{\frac{\partial #1}{\partial #2}}
\newcommand{\doublePartDeriv}[2]{\frac{\partial^2 #1}{\partial^2 #2}}


\title{Problem 2.6}
\author{Gabriel Decker}


\begin{document}


\maketitle
\begin{flushleft}


In this problem, we find the electric field a distance $z$ above the center of a circular disk of charge with charge density $\sigma$ and radius $R$.
For this problem, define the coordinate system such that the center of the disk is at the origin, the ring exists in the $xy$-plane, and the $\hat{z}$ direction points from the origin to the point of interest.
Since we are dealing with a surface of charge, will use the following equation to find our electric field.
$$\vec{E}(\vec{r})=\frac{1}{4\pi\epsilon_0}\int\frac{\sigma(\vec{r'})}{\scripty{r}^2}\hat{\scripty{r}}\,da'$$\\~\\

We will use the following result from problem 2.5, where we found the electric field a distance $z$ above the center of a ring of line charge with radius $r$.
$$\vec{E}(z,r)=\frac{1}{4\pi\epsilon_0}\frac{r\lambda}{(r^2+z^2)^{3/2}}\int_0^{2\pi}-r\hat{s}+z\hat{z}\,d\phi$$
$$\implies\vec{E}(z,r)=\frac{1}{4\pi\epsilon_0}\frac{2\pi r\lambda z}{(r^2+z^2)^{3/2}}\hat{z}$$\\~\\

You can safely make the assumption that we will treat our disk as a bunch of rings and sum up the total electric field for all the infinitesimal rings with radii between zero and $R$, but I will demonstrate that this assumption works.\\~\\

Let's start by finding how $\lambda$ is related to $\sigma$.
$\sigma$ is total charge $Q$ divided by the total surface area $\pi R^2$, so $\sigma=Q/\pi R^2$.
One ring of the disk at radius $r$ will take a sliver of that area.
It has area $\pi(r+dr)^2-\pi r^2=\pi(r^2+2rdr+dr^2)-\pi r^2=2\pi rdr$ (neglecting the doubly infinitesimal $dr^2$).


\end{flushleft}

\end{document}
